\section{A Calibrated Synthetic ERP Economy}\label{sec:synthetic}

The four real datasets used elsewhere in this paper give us, at best, a few
thousand labeled fraud rows each, concentrated in typologies each simulator's
original authors happened to encode. To probe typologies we do not observe at
scale in any single real ledger, and to obtain a dataset that is simultaneously
account-to-account (so that cycles are structurally possible), ERP-flavored
(payables, payroll, receivables), and large enough in fraud count to support
per-typology analysis, we built \texttt{synth\_erp}: an agent-based simulation
of a 40-company economy running for 365 simulated days over $\sim$39k accounts,
producing 1{,}148{,}503 transactions of which 5{,}609 are labeled fraud
(a 0.4884\% rate). This section documents how the generator was designed to
resist the one failure mode that would make it useless for our purposes, what
an independent realism audit against the four real datasets found — including
where the synthetic economy does \emph{not} match reality — and what a direct
adversarial probe for cheap tabular leaks and a downstream training-transfer
experiment revealed. Throughout, we hold \texttt{synth\_erp} to the same
evidentiary standard as the real datasets: numbers are reported whether or not
they are flattering, because the entire purpose of this exercise is to know
which sanctioned use cases actually survive scrutiny.

\subsection{Design Philosophy: The Circularity Trap}
\label{sec:synthetic-trap}

A synthetic fraud generator built by the same team that builds the fraud
\emph{detector} is exposed to an obvious failure mode: if the generator's
authors hand-encode fraud as ``insert a 3-hop cycle'' or ``make the receiver's
in-degree spike,'' a topology engine built by the same authors will of course
find it — not because graph structure is informative about fraud in general,
but because the two pieces of code were written to fit each other. Grading a
detector against a generator that shares its assumptions proves nothing about
the real world; it is closer to checking that a lock opens with its own key.
\texttt{synth\_erp} was designed against this failure mode by a single binding
rule, carried in this project's own methodology documentation: fraud is
injected as \emph{goal-driven behavior with evasion}, never as a
detector-shaped pattern. The fraud-generation code imports nothing from the
topology-feature codebase and references no detector's internals; every
typology is written as an agent pursuing a criminal objective (skim a
receivable, drain a payables account, collect from recruited mules) under
realistic constraints (approval thresholds, banking hours, evasion timing),
and any cycle, fan-in, or amount-recurrence a detector later measures is a
\emph{consequence} of that behavior, not its specification. The benign
population mirrors this discipline: ordinary commerce is allowed to produce
the same raw ingredients fraud produces — busy retailers with heavy fan-in,
payroll batches that resemble fan-out bursts, subscription billing that
recurs at a near-equal amount, refunds that close reciprocal pairs, multi-hop
circulation as salaries fund purchases that fund vendors — since a background
in which only fraud has structure would reintroduce the same shortcut through
the back door. Sec.~\ref{sec:synthetic-emergence} reports just how modestly
enriched the fraud-typology structural traces actually are once this
constraint is honored — a dataset engineered to make cycles trivially
separable would show the opposite pattern.

\subsection{Generator Architecture, Scale, and Transaction Composition}

\texttt{synth\_erp} follows the same broad methodology family as two of the
real datasets used elsewhere in this paper — an agent-based simulation
calibrated against measured real-world statistics rather than a purely
hand-tuned generator, the approach behind BankSim~\cite{lopezrojas2014banksim},
PaySim~\cite{lopezrojas2016paysim}, and the AMLSim lineage behind
\texttt{ibm\_aml}~\cite{altman2023realistic}. The simulated economy comprises
40 companies operating alongside their customers, vendors, and employees,
generating seven ordinary transaction flows
(\texttt{sale}, \texttt{vendor\_payment}, \texttt{payroll}, \texttt{transfer},
\texttt{refund}, \texttt{reimbursement}, \texttt{recurring\_bill}) over a full
simulated year. Table~\ref{tab:synth-flows} gives the exact per-flow
transaction counts, read directly from the same adversarial tabular-leak audit
discussed in Sec.~\ref{sec:synthetic-leakprobe} (the per-\texttt{tx\_type}
support counts there sum exactly to the dataset total, confirming the two are
consistent measurements of the same underlying table). Retail \texttt{sale}
traffic dominates by volume (66.2\%), with the ERP-specific payables rail
(\texttt{vendor\_payment}, 18.1\%) the second-largest flow — a composition
deliberately closer to an operating company's ledger than to a pure
consumer-payments simulator.

\begin{table}[t]
\centering
\small
\caption{Transaction volume by flow type, \texttt{synth\_erp} (all 1{,}148{,}503 rows, fraud and benign combined). Source: \texttt{paper/data/synth\_erp\_realism.json}, \texttt{tabular\_leak\_probe.rules} per-\texttt{tx\_type} support.}
\label{tab:synth-flows}
\begin{tabular}{@{}lrr@{}}
\toprule
Flow (\texttt{tx\_type}) & Count & Share \\
\midrule
\texttt{sale}            & 760{,}558 & 66.23\% \\
\texttt{vendor\_payment}  & 208{,}155 & 18.13\% \\
\texttt{transfer}         &  94{,}781 &  8.25\% \\
\texttt{payroll}          &  49{,}240 &  4.29\% \\
\texttt{refund}           &  18{,}396 &  1.60\% \\
\texttt{reimbursement}    &  12{,}236 &  1.07\% \\
\texttt{recurring\_bill}  &   5{,}137 &  0.45\% \\
\midrule
Total                     & 1{,}148{,}503 & 100.00\% \\
\bottomrule
\end{tabular}
\end{table}

At the account/graph level, the simulated economy is dense and repeat-heavy:
39{,}256 accounts transact across 251{,}192 unique sender--receiver pairs,
with 85.86\% of transactions occurring between a pair that has transacted
before (\texttt{repeat\_pair\_txn\_frac} = 0.8586) and 94.1\% of accounts
appearing on both the sending and receiving side of at least one transaction
(\texttt{namespace\_overlap} = 0.941) — i.e., this is deliberately \emph{not}
a bipartite consumer-to-merchant graph, since ERP economies circulate money
(salaries fund purchases, purchases fund vendors, vendors pay employees) in a
way that pure-payments simulators structurally cannot. Degree is heavy-tailed
in both directions: median out-degree 18 (p99 = 83, max = 20{,}838) and
median in-degree 4 (p99 $\approx$ 128, max = 194{,}265), driven by a small
number of high-volume retailers and payroll-issuing companies. Activity is
temporally realistic without being explicitly targeted at any metric: 79.06\%
of transactions fall in business hours and 17.31\% on weekends, at a mean of
3{,}146.6 transactions/day (p99 = 5{,}308.3) sustained across all 365 simulated
days — timestamped to the second, a finer time axis than any of the four real
datasets in our roster provide.

\subsection{Five Fraud Typologies as Emergent Behavior}
\label{sec:synthetic-emergence}

Fraud is injected as five occupational-fraud typologies, each defined by an
actor's goal and evasion strategy rather than by a target graph shape:
\texttt{laundering\_cycle} (embezzled funds layered through a ring of shells
and, in most scenarios, cycled back toward the entry account), \texttt{mule\_fanin}
(scam proceeds collected from many recruited mule accounts into one
collector), \texttt{lapping} (a diverted receivable covered days later by the
next customer's payment), \texttt{invoice\_kickback} (a colluding vendor
overbills and returns a cut to an insider), and \texttt{larceny} (direct
insider diversion of funds, disguised as vendor or reimbursement traffic). The
resulting typology counts are 1{,}442 (\texttt{laundering\_cycle}), 2{,}327
(\texttt{mule\_fanin}), 760 (\texttt{lapping}), 625 (\texttt{invoice\_kickback}),
and 455 (\texttt{larceny}) labeled rows, summing to the reported 5{,}609
fraud transactions.

Because these behaviors are never specified in terms of detector features, the
honest question is whether they leave structural traces at all, and at what
rate. The benign background already contains a nontrivial rate of graph
cycles — 5.761\% of benign transactions close a cycle in the observed window,
with 99th-percentile benign fan-in reaching 4{,}197 distinct senders (busy
retailers and payroll hubs) — so cyclicity and fan-in are not, by
construction, fraud-exclusive signals. Against that background,
\texttt{laundering\_cycle} transactions close a cycle 18.45\% of the time (all
at length $\geq 3$), a $3.2\times$ enrichment over the 5.761\% benign base
rate — real, but far from deterministic, consistent with a design in which
60\% of rings loop back to their entry account and 40\% exit linearly, and in
which the streaming, look-ahead-free topology engine's window will miss some
closures entirely. \texttt{mule\_fanin} collectors show a 90th-percentile
in-window fan-in of only 8 distinct senders — far below the benign
99th-percentile of 4{,}197 — so a raw fan-in threshold offers no free
discrimination there either. \texttt{lapping} shows the weakest structural
footprint of the five: its amount-recurrence measure
(\texttt{lap\_in\_recur}) averages 0.06, swamped by a benign population
dominated by genuinely recurring payroll and subscription billing. We flag
this pattern explicitly because it foreshadows a headline finding developed
later in this paper (Sec.~\ref{sec:ablation}): engineered cycle-shaped
features carry a surprisingly small share of the detectable signal on real
data, and this section shows that is not an artifact of an unrealistic
generator — even in a dataset whose typologies are explicitly designed to
produce cycles and fan-in, those structural traces are only modestly
enriched over an already-active benign background, not deterministic
markers.

\subsection{Realism Audit Against the Real-Data Reference Band}
\label{sec:synthetic-realism}

For every statistic we report on \texttt{synth\_erp}, we ask one question:
does it fall inside the range spanned by the four real datasets
(\texttt{ibm\_aml}, \texttt{banksim}, \texttt{sap\_wurzburg}, \texttt{paysim}),
measured with identical methodology, or outside it? Table~\ref{tab:synth-realism}
answers this for six statistics drawn from \texttt{calibration\_reference.json}
and \texttt{synth\_erp\_realism.json}. The fraud rate, the repeat-pair share
of transactions, and the sender/receiver namespace overlap all land inside
the real span. Three do not, and we report them plainly rather than tuning
them away. Log-amount dispersion (\texttt{log1p\_std} = 2.3256) sits above
every real dataset (span 0.9551--2.0193): the simulated economy mixes
cents-scale retail tickets with treasury sweeps of hundreds of thousands of
dollars in one distribution, a wider mixture than any single-domain real
dataset exhibits — genuine \emph{over}-dispersion, not a defensible match.
Benford mean-absolute-deviation (\texttt{benford\_mad} = 0.0047) is
\emph{lower} than every real dataset (span 0.0142--0.0758): \texttt{synth\_erp}'s
amounts obey Benford's law more cleanly than any real ledger we measured,
itself a realism failure in the opposite direction — real accounting data
carries idiosyncratic roundness and human-entered biases that a mixture of
lognormal generators does not reproduce. Pair reciprocity
(\texttt{pair\_reciprocity} = 0.1356) is roughly $3.0\times$ the real maximum
(0.0446, \texttt{ibm\_aml}): salary-to-purchase circulation and refund/cover
behavior generate more back-and-forth pairs than any of the four real
economies structurally permit. None of these three deviations is hidden in
downstream use: they are why \texttt{synth\_erp} is never treated as evidence
for this paper's central topology claim, only as training material and a
stress test, per the usage policy at the top of this section.

%% FIGURE TODO: grouped bar chart with one bar-triplet per calibration metric
%% (fraud rate, log1p_std, benford_mad, pair_reciprocity, repeat_pair_frac,
%% namespace_overlap): synth_erp value vs. the real-dataset min-max band
%% (shown as a shaded range or error bar), so in-range vs. out-of-range
%% metrics are visually obvious at a glance. Omitted from the body to save
%% space given Table~\ref{tab:synth-realism} already reports the same numbers.

\begin{table}[t]
\centering
\small
\caption{Realism audit: \texttt{synth\_erp} vs.\ the real-dataset reference band (min--max across \texttt{ibm\_aml}, \texttt{banksim}, \texttt{sap\_wurzburg}, \texttt{paysim}). Source: \texttt{calibration\_reference.json}, \texttt{synth\_erp\_realism.json}.}
\label{tab:synth-realism}
\resizebox{\columnwidth}{!}{%
\begin{tabular}{@{}lrrrl@{}}
\toprule
Metric & \texttt{synth\_erp} & Real min & Real max & In band? \\
\midrule
Fraud rate               & 0.4884\% & 0.1236\% & 1.2108\% & Yes \\
Repeat-pair txn.\ share  & 0.8586   & 0.0      & 0.9936   & Yes \\
Namespace overlap        & 0.9410   & 0.0      & 0.9927   & Yes \\
log1p(amount) std.\      & 2.3256   & 0.9551   & 2.0193   & \textbf{No (above)} \\
Benford MAD              & 0.0047   & 0.0142   & 0.0758   & \textbf{No (below)} \\
Pair reciprocity         & 0.1356   & 0.0      & 0.0446   & \textbf{No (3.0$\times$ above)} \\
\bottomrule
\end{tabular}%
}
\end{table}

A more direct realism concern is whether fraud amounts alone give fraud away —
the specific artifact this project identified in \texttt{ibm\_aml} elsewhere in
this paper (fraud amounts distinctively \emph{small}, amount-ratio 0.068,
globally separable at ROC-AUC $\approx 0.91$). On \texttt{synth\_erp}, the
median fraud amount (\$1{,}692.59) is 23.79$\times$ the median benign amount
(\$71.14; the equivalent ratio in \texttt{banksim}, from
\texttt{calibration\_reference.json}, is 11.9$\times$), and a univariate
log-amount classifier reaches ROC-AUC 0.8078 overall (KS statistic 0.604
between the fraud and benign log-amount distributions) — at first glance a
fraud population skewed to B2B-sized amounts in a retail-heavy economy.
Restricted to the amount band where 93\% of fraud actually lives
(\$300--\$20{,}000, $n=5{,}228$ fraud rows in band), that same classifier
collapses to ROC-AUC 0.4911 — indistinguishable from chance. Fraud amounts
are globally distinguishable only because fraud skews toward a coarser amount
regime than the bulk of retail traffic; \emph{within} its own regime, amount
carries essentially no signal, forcing any detector toward behavioral or
structural evidence rather than a one-column amount rule — the opposite of
the \texttt{ibm\_aml} artifact, and a deliberate design target corroborated
here after the fact rather than asserted in advance.

\subsection{Adversarial Tabular-Leak Probe}
\label{sec:synthetic-leakprobe}

A realistic amount distribution does not rule out a cheap categorical
shortcut elsewhere in the table — the same class of hidden bias that motivates
auditing tabular fraud-evaluation datasets for shortcut features before
trusting them~\cite{jesus2022turning}. We therefore ran an automated probe over 16
candidate rules — every \texttt{tx\_type} indicator, an account-creation-date
indicator on each side of the transaction, and pairwise conjunctions of the
two — scoring each rule's support, precision, recall of fraud, and lift over
the 0.4884\% base rate. Table~\ref{tab:synth-leakprobe} reports the eight
highest-lift rules out of the 16 tested. Three flow types
(\texttt{tx\_type=payroll}, \texttt{tx\_type=recurring\_bill},
\texttt{tx\_type=refund}) and their conjunctions with a mid-year account
carry \emph{zero} fraud whatsoever (lift 0.0) — no scenario in the generator
routes any of the five typologies onto these three rails at all, which is a
genuine (if narrow) gap in coverage rather than a realism strength.

The rule the probe itself records as worst,
\texttt{tx\_type=transfer \& dest\_created\_midyear}, reaches a lift of
$19.0\times$ the base rate, support of 26{,}420 transactions, precision of
9.29\%, and recovers 43.75\% of all fraud in a single two-clause rule. The
dataset's own automated gate does not classify this as a leak
(\texttt{leak\_flag} = \texttt{false}, its precision bar sitting well above
what this rule reaches), and a rule flagging under 10\% of the transactions it
fires on as fraud is far from the near-deterministic, one-line shortcuts the
generator's earlier drafts contained before an adversarial review removed
them. Still, a rule with 19$\times$ lift and 44\% recall is a real,
exploitable regularity, not a non-finding, and we report it rather than let a
binary pass/fail flag stand in for a reader's own judgment of where
``legitimate risk factor'' ends and ``tabular leak'' begins. A transaction
moving through the \texttt{transfer} rail toward an account opened within the
same year is a defensible real-world risk heuristic — mid-year account
creation correlates with fraud in practice, not only in this simulation — but
we flag it here precisely so that any model trained on \texttt{synth\_erp}
can be checked for whether it quietly learned this shortcut instead of the
intended structural and behavioral signal.

\begin{table}[t]
\centering
\scriptsize
\caption{Adversarial tabular-leak probe, top 8 of 16 rules by lift. Base fraud rate 0.4884\%. Source: \texttt{synth\_erp\_realism.json}, \texttt{tabular\_leak\_probe}.}
\label{tab:synth-leakprobe}
\resizebox{\columnwidth}{!}{%
\begin{tabular}{@{}lrrrr@{}}
\toprule
Rule & Support & Precision & Recall & Lift \\
\midrule
\texttt{transfer \& dest\_midyear}$^\dagger$ & 26{,}420 & 0.0929 & 0.4375 & \textbf{19.0} \\
\texttt{dest\_created\_midyear}         & 67{,}621  & 0.0502 & 0.6049 & 10.3 \\
\texttt{tx\_type=transfer}              & 94{,}781  & 0.0385 & 0.6500 & 7.9 \\
\texttt{vendor\_payment \& dest\_midyear}$^\dagger$ & 25{,}781 & 0.0316 & 0.1451 & 6.5 \\
\texttt{sale \& dest\_midyear}$^\dagger$          & 8{,}131   & 0.0146 & 0.0212 & 3.0 \\
\texttt{src\_created\_midyear}          & 156{,}550 & 0.0113 & 0.3147 & 2.3 \\
\texttt{reimbursement \& dest\_midyear}$^\dagger$ & 741       & 0.0081 & 0.0011 & 1.7 \\
\texttt{tx\_type=reimbursement}         & 12{,}236  & 0.0079 & 0.0173 & 1.6 \\
\bottomrule
\end{tabular}%
}
\vspace{2pt}
\raggedright\scriptsize $^\dagger$\texttt{dest\_midyear} abbreviates \texttt{dest\_created\_midyear}.
\end{table}

\subsection{The Augmentation Transfer Probe}
\label{sec:synthetic-augmentation}

The realism and leak audits above establish what \texttt{synth\_erp} is
calibrated well enough to be used \emph{for}. The most direct test of the
narrower, practical question — does adding \texttt{synth\_erp} rows to a real
dataset's training split make its fraud model better? — is a paired
augmentation experiment: for each of three real datasets, train an XGBoost
model~\cite{chen2016xgboost} three ways (real data only; real data with all
1{,}148{,}503 \texttt{synth\_erp} rows pooled into the training split; and
\texttt{synth\_erp} data only, evaluated zero-shot on the real test split),
each averaged over three seeds, under the same time-aware split, class-weight
imbalance handling, and validation-tuned threshold used throughout this paper.
Table~\ref{tab:synth-augmentation} reports the result, and it is unambiguous:
pooling synthetic rows into training \emph{hurts} test PR-AUC on every one of
the three real datasets tested. \texttt{banksim} falls from
$0.8669\pm0.0031$ (real-only) to $0.7610\pm0.0022$ (real+synth), a paired
per-seed delta of $-0.1059\pm0.0011$; \texttt{ibm\_aml} falls from an
already-saturated $0.99999\pm0.00001$ to $0.9954\pm0.0013$,
$\Delta=-0.0046\pm0.0013$ — small in absolute terms, but the baseline leaves
almost no headroom to fall into, so even this small drop is a consistent,
same-direction result; and \texttt{sap\_wurzburg} — on only 25 test-set
frauds, the same small-sample caveat that applies to every
\texttt{sap\_wurzburg} result in this paper — falls from $0.1366\pm0.0446$ to
$0.0293\pm0.0033$, $\Delta=-0.1073\pm0.0417$.
The mechanism is unsurprising in hindsight given Sec.~\ref{sec:synthetic-realism}:
a single tree ensemble fit jointly over two economies with measurably
different amount dispersion and reciprocity structure pulls its splits toward
the pooled distribution, at the real dataset's expense. We report this as a
clean negative result rather than qualify it away: \emph{naive} pooled
augmentation with this dataset, under this protocol, does not help, on any
real dataset we tested it against.

\begin{table}[t]
\centering
\small
\caption{Augmentation transfer probe: XGBoost PR-AUC, 3-seed mean $\pm$ std. \texttt{synth\_only} is zero-shot (trained solely on \texttt{synth\_erp}, evaluated on the real test split). Source: \texttt{augmentation.json}.}
\label{tab:synth-augmentation}
\resizebox{\columnwidth}{!}{%
\begin{tabular}{@{}lrrrr@{}}
\toprule
Dataset & Real-only & Real+synth & $\Delta$ & Synth-only \\
\midrule
\texttt{banksim} ($n{=}1320$)      & 0.8669$\pm$0.0031 & 0.7610$\pm$0.0022 & $-$0.1059$\pm$0.0011 & \textbf{0.7243$\pm$0.0079} \\
\texttt{ibm\_aml} ($n{=}377$)      & 0.99999$\pm$0.00001 & 0.9954$\pm$0.0013 & $-$0.0046$\pm$0.0013 & 0.0017$\pm$0.0006 \\
\texttt{sap\_wurzburg} ($n{=}25$)  & 0.1366$\pm$0.0446 & 0.0293$\pm$0.0033 & $-$0.1073$\pm$0.0417 & 0.0008$\pm$0.0001 \\
\bottomrule
\end{tabular}%
}
\end{table}

There is exactly one genuinely positive number in this experiment, and it is
a narrower, different claim than ``synthetic data helps supervised
training'': the \texttt{synth\_only} zero-shot column for \texttt{banksim}.
A model that has \emph{never seen a single real BankSim row} — trained
entirely on \texttt{synth\_erp}, evaluated as-is on real, held-out
\texttt{banksim} test data — scores $0.7243\pm0.0079$ PR-AUC, recovering
83.6\% of the $0.8669$ real-trained baseline, at precision@100 of
$0.99\pm0.0082$ against the real-trained baseline's perfect $1.00$; chance at
this split's 1.21\% fraud rate is on the order of $0.01$--$0.02$ PR-AUC, so
this is far from a vacuous transfer. The same zero-shot comparison fails on
\texttt{ibm\_aml} ($0.0017$ PR-AUC against a $0.99999$ real-trained ceiling —
unsurprising, since \texttt{ibm\_aml}'s giveaway-small fraud amounts are
inverted relative to \texttt{synth\_erp}'s B2B-skewed fraud) and on
\texttt{sap\_wurzburg} ($0.0008$ PR-AUC, 25 test frauds, too small to support
any transfer claim). We read the \texttt{banksim} result as evidence that the
synthetic economy teaches \emph{transferable structural and behavioral
discrimination} that generalizes to a real domain it structurally resembles —
not as evidence that pooling synthetic rows into a real training set is a
useful augmentation recipe. These are two different claims, only one of which
this experiment supports; we are careful not to let the positive framing of
the zero-shot result launder the negative framing of the pooled-augmentation
result three rows above it.

\subsection{Summary and Scope}

\texttt{synth\_erp} is a large, ERP-flavored, multi-typology synthetic ledger
built under a binding non-circularity constraint and audited, after the fact,
against four real datasets on realism, tabular leakage, and downstream
transfer. It matches the real fraud-rate, repeat-pair, and namespace-overlap
bands; it is measurably over-dispersed in amount and reciprocity and
unrealistically Benford-clean relative to every real dataset compared against
it; its cheapest categorical shortcut reaches 19$\times$ lift and 44\% recall
without being flagged as a formal leak, a judgment call we surface rather
than hide; and pooling it into real training data hurts every real dataset's
PR-AUC we tested, its one positive contribution being a narrower zero-shot
transfer result on the one real dataset it structurally resembles.
Consistent with the usage policy stated at the outset, none of this
section's numbers are used anywhere in this paper as evidence for the
central topology claim (Sec.~\ref{sec:ablation}); \texttt{synth\_erp}'s role
is training material, a per-typology stress test, and demo content for the
audit console (Sec.~\ref{sec:system}), never proof.
